% ====================================================================
% ch1_sec05b_gr2L.tex — [GR-2L] stage-2L 의 엔트로피 온도 분리 + 두-상/고용체 판정자
%   삽입: §5 폭·이중지위(sec:width) 뒤. v1.0.24 신규(@5 반영).
%   설계 정합(master 재조정): 본 소절은 §7(sec:broadening-class)의 두-상/고용체
%     분류를 뒤집지 않는다 — 판정자와 stage-2L 온도 서명(핵심 @5)을 얹되, 전이별
%     상성격 확정은 §7 + 피팅된 Ω 로 위임(P5·gate-6·gate-7). 5-feature 세분은
%     '추가 후보'(선택적 코드 시드)로만 제시, tab:staging 4-전이 기준은 보존.
%   refine_b(R1 Phase-②): 분류 위임을 문장 단위로 airtight 화(판정자=기준식이지 판결이
%     아님을 도입부·(a)·(d) 에서 일관 명시) + 5-feature↔4-전이 명칭대응을 '4 + 추가 1'
%     로 정정(brief 의 '두-상4+고용체1' 재배치 잔재 제거). 강점(판정자 유도·boxed
%     U(T) 지도·schmitt2022·−0.44%p·w_eff 다리·warnbox)은 전부 보존.
% ====================================================================
\subsection{stage-2L 의 엔트로피 온도 분리 --- 판정자 $\dd\mu/\dd\theta|_{1/2}=4RT-2\Omega$ 와 그 온도 서명}\label{ssec:gr-2L}
앞 소절(\S\ref{sec:width})이 폭 $w_j=n_jRT/F$(식~\eqref{eq:wbase})의 \emph{이중지위} --- 단상($\Omega_j\le2RT$)에서는
평형 예측, 두-상($\Omega_j>2RT$)에서는 broadening 이 정하는 현상학적 폭 --- 를 세웠고, 그 지위를 가르는 것이
전이별 $\Omega_j$ 임을 밝혔다. 이 소절은 그 위에 둘을 얹는다: (i) 어떤 전이가 두-상이고 어떤 것이 고용체인지의
\emph{판정자}(전이별 분류를 \emph{확정}하는 도구가 아니라 그 판정 \emph{기준} --- 확정은 \S\ref{sec:broadening-class}
소관)를 정칙용액 화학퍼텐셜에서 닫고, (ii) 흑연 staging 의 가운데 쌍($3\!\to\!2\mathrm L,\,2\mathrm L\!\to\!2$)이
실은 \emph{엔트로피 안정화 중간상}(stage-2L)을 사이에 둔 분리쌍임을 온도 서명으로 보인다. 새 물리는 없다 ---
식~\eqref{eq:gpp}$\cdot$\eqref{eq:spinodal}$\cdot$\eqref{eq:Uj} 의 흑연 골격을 읽을 뿐이다.
\textbf{경계(가드)} --- 본 소절은 전이별 두-상/고용체 \emph{분류}를 새로 확정하지 않는다: 그 분류는
\S\ref{sec:broadening-class} 의 실측-plateau 기준을 그대로 따르며(아래 검산 참조), 여기서 더하는 것은
\emph{판정자}(분류의 판정식)와 \emph{stage-2L 온도 분리}(가운데 쌍의 $\Delta S$ 차 서명)뿐이다.
\textbf{역사적 명칭 가드} --- stage 번호($1\cdot2\cdot2\mathrm L\cdot3\cdot4$)는 Dahn 상도표\cite{dahn1991} 의
\emph{구조상} 명명이며, ``ver.$N$''(문건 이력)$\cdot$``Chapter $N$''(문서 구조)과 무관하다.

\textbf{(a) 판정자 --- $\dd\mu/\dd\theta|_{1/2}=4RT-2\Omega$.} 어떤 전이가 두-상 plateau 로 갈 수 있는지 연속
고용체로 남는지를 \emph{가르는 기준}은 정칙용액 자유에너지 식~\eqref{eq:gxi} 의 화학퍼텐셜 부호가 준다(이 기준이
전이별 분류를 혼자 \emph{확정}하지는 않고 \S\ref{sec:broadening-class} 의 실측 plateau 와 맞물리는 방식은 아래
검산에서 닫는다). 전이 창 안 자리 점유 $\theta$(리튬화 진행)에
대한 자리 화학퍼텐셜은 평균장 정칙용액에서
\begin{equation}
\mu(\theta)=\mu^\circ+RT\ln\frac{\theta}{1-\theta}+\Omega\,(1-2\theta)
\label{eq:gr2l-mu}
\end{equation}
이고($\mu=\partial g/\partial\theta$ 를 $\theta$-좌표로 적은 것), 이를 한 번 더 미분하면 --- 로그 항은
$\dd/\dd\theta\,\ln[\theta/(1-\theta)]=1/\theta+1/(1-\theta)=1/[\theta(1-\theta)]$, 상호작용 항은
$\dd/\dd\theta\,\Omega(1-2\theta)=-2\Omega$ --- 등온선의 국소 강성이 나온다:
\begin{equation}
\frac{\dd\mu}{\dd\theta}=\frac{RT}{\theta(1-\theta)}-2\Omega
\quad\Longrightarrow\quad
\frac{\dd\mu}{\dd\theta}\Big|_{\theta=1/2}=\frac{RT}{1/4}-2\Omega=4RT-2\Omega.
\label{eq:gr2l-disc}
\end{equation}
이는 흑연 곡률식~\eqref{eq:gpp}($\partial^2 g/\partial\xi^2$)를 대칭점에서 읽은 것과 \emph{같은 대수}다
($\theta{=}1{-}\xi$ 대칭에서 $\theta{=}1/2\leftrightarrow\xi{=}1/2$). 판정 기준은 이 한 부호다: $\Omega>2RT$ 면
$\dd\mu/\dd\theta|_{1/2}<0$ 라 등온선이 비단조(back-bending)이고 spinodal gap(식~\eqref{eq:spinodal})이 열려
\emph{miscibility gap$\cdot$두-상 공존}(평형 등온선이 near-delta)이며, $\Omega<2RT$ 면 등온선이 단조라 gap 이 없어
\emph{연속 고용체}(유한 폭 종)다. $2RT\approx4958$ J/mol$@298.15$ K 가 문턱이다(식~\eqref{eq:spinodal} 와 동일).
곧 \S\ref{sec:width} 이중지위의 두 갈래는 임의 분류가 아니라 이 한 부등호 $\Omega_j\lessgtr2RT$ 의 결과다(다만 이
부등호가 실재 전이를 두-상/고용체로 \emph{확정}하는 것은 아님은 바로 아래 검산에서 \S\ref{sec:broadening-class}
과 맞춘다).

\textbf{단일상 쪽 $w_\eff$ 가 재는 것 --- 높이는 정확, 폭은 아니다(v1.0.25 정정).} 단일상($\Omega<2RT$)
등온선의 국소 강성에서 나오는 척도
\[
w_\eff\;\equiv\;\frac{RT}{F}\Big(1-\frac{\Omega}{2RT}\Big)\;=\;\frac{RT}{F}\,\lambda,
\qquad \lambda\equiv1-\frac{\Omega}{2RT}>0\ \ (\Leftrightarrow\ \Omega<2RT),
\]
는 판정자~\eqref{eq:gr2l-disc} 의 $\theta{=}\tfrac12$ 강성에 반비례하며, 이때 \emph{중심 높이}는
$(\dd Q/\dd V)|_{\theta=1/2}=QF/(4RT-2\Omega)=Q/(4w_\eff)$ 로 \emph{정확}하다(모든 $\lambda$ 에서 항등 ---
수치 확인). 그러나 이 $w_\eff$ 를 \emph{반높이 폭}으로 읽어서는 안 된다. 반높이 조건
$\theta(1-\theta)=1/[4(1+\lambda)]$ 를 풀면 $\theta=\tfrac12\pm\tfrac12x$, $x\equiv\sqrt{\lambda/(1+\lambda)}$
이고, 등온선 $V(\theta)$ 에 넣으면 반높이 폭이 닫힌 형태로 나온다:
\begin{equation}
\mathrm{FWHM}=\frac{4RT}{F}\,\mathrm{artanh}\,x-\frac{2\Omega}{F}\,x
\;\;\xrightarrow[\;\lambda\to0\;]{}\;\;\frac{16}{3}\frac{RT}{F}\,\lambda^{3/2}
\qquad\Big(x=\sqrt{\tfrac{\lambda}{1+\lambda}}\Big).
\label{eq:gr2l-fwhm}
\end{equation}
곧 폭은 $\lambda$ 가 아니라 $\lambda^{3/2}$ 로 닫힌다 --- 까닭은 $u\equiv\theta-\tfrac12$ 전개
$V(\theta)-U^\circ=-\tfrac{4RT}{F}\lambda\,u-\tfrac{16RT}{3F}u^3+\mathcal O(u^5)$ 에서 \emph{1차항 계수만}
$\lambda$ 에 비례해 소멸하고 3차항(자유에너지 4차항)은 $\lambda$ 와 무관하게 남기 때문이다. 그러므로 $w_\eff$ 를 폭으로
읽는 순진한 대입은 $\lambda\to0$ 에서 실제 폭을 $\approx0.66/\sqrt{\lambda}$ 배 \emph{과대}평가한다 ---
$\lambda=10^{-2}$ 에서 약 $6.6$ 배, $10^{-3}$ 에서 약 $21$ 배, $5\times10^{-4}$ 에서 약 $30$ 배다(수치 확인).
점근식~\eqref{eq:gr2l-fwhm} 자체의 정확도도 함께 적어 둔다 --- 좌변 닫힌형 대비 과대율이
$\lambda=0.01$ 에서 $0.6\%$, $0.02$ 에서 $1.2\%$, $0.05$ 에서 $3.0\%$, $0.5$ 에서 $27\%$ 이므로, $\lambda\gtrsim0.02$
에서는 점근식이 아니라 \emph{좌변 닫힌형}을 쓴다. 따라서 이 식은 두-상 극한으로의 \emph{연속 보간이 아니다} --- 단일상 쪽 \emph{중심
높이}의 국소 지표일 뿐이고, 판정자가 $\Omega=2RT$ 를 넘는 순간 평형 등온선은 Maxwell 공존평탄으로
\emph{불연속}하게 갈아탄다(near-delta). 두-상($\Omega>2RT$)의 \emph{유한} 폭은 평형이 예측하지 않으며
\S\ref{sec:broadening} 의 세 출처$\cdot$현상학적 피팅이 정한다 --- 폭식~\eqref{eq:wbase} 의 이중지위는
이 지점에서 \emph{두 지위로 갈리는} 것이고 한 식으로 이어지는 것이 아니다.

\begin{verifybox}
검산 --- \emph{판정자와 §7 분류의 정합(과대주장 금지).} 정칙용액 피팅 진단(regsol2)이 표~\ref{tab:staging}
네 전이의 $\Omega_j/RT$ 를 $[\,4.06,\,2.02,\,3.55,\,4.07\,]$ 로 되돌려 전부 $2$ 를 넘는다. 이는 그 자체로 네
전이를 모두 두-상으로 \emph{확정}하지 않는다 --- \S\ref{sec:broadening-class} 가 이미 밝혔듯 표의 \emph{초기}
$\Omega_j$ 로는 네 전이가 모두 문턱을 넘으므로, 실제 두-상/고용체 지목의 근거는 $\Omega$ 문턱이 아니라
\emph{실측 plateau$\cdot$staging 상도표}\cite{dahn1991,ohzuku1993} 다. 따라서 본 소절의 판정자~\eqref{eq:gr2l-disc}
는 분류를 \emph{다시 여는} 도구가 아니라, \S\ref{sec:broadening-class} 의 분류가 어느 부등호에서 나오는지를
닫는 판정\emph{식}이며, 최종 두-상 확정은 \emph{피팅된} $\Omega_j$ 의 $2RT$ 초과 $+$ 실측 plateau 다. 경계값
$2.02$($3\!\to\!2\mathrm L$ 슬롯)는 문턱 바로 위라 시료$\cdot$도핑 편차로 재분류될 수 있는 민감 지점으로 남긴다.
\end{verifybox}

\begin{srcbox}
\textbf{다리 --- 단일 $\Omega_j$ 는 다중-부분격자 staging 질서의 평균장 축약이다.} 식~\eqref{eq:gr2l-mu} 의 한
모수 $\Omega_j$ 는 흑연 staging 의 서로 다른 두 상호작용 --- 면내(in-plane) 인력과 층간(staging) 질서 --- 을
한 진행좌표로 접은 것이다. 제일원리 계산\cite{persson2010b} 은 이 다층 질서(Daumas--H\'erold 교대충전)가
stage 1$\cdot$2 의 sharp 두-상은 강하게 재현하되 stage 3$\cdot$4 영역은 신뢰 밖으로 두는 구조임을 보인다 ---
곧 실제 흑연은 단일 $\Omega$ 가 아니라 다중-부분격자 ordering 이다. 본 소절의 단일 평균장 $\Omega_j$(정규용액
이중웰)는 그 다체 구조 중 \emph{면내 응축의 sharpness 만} 유효 문턱으로 포착하고 층간 다층 질서는 담지 못하는
근사이며, 판정자~\eqref{eq:gr2l-disc} 는 그 근사 \emph{안에서만} 두-상/고용체의 유효 문턱을 정한다(전이별
면내$\cdot$층간 $\Omega$ 분해는 회사 피팅 단계로 위임, 본 소절은 유효 단일 $\Omega_j$ 만 확정).
\end{srcbox}

\textbf{(b) stage-2L 은 엔트로피 안정화 상 --- 그 온도 서명은 $\Delta(\Delta S)$.} 가운데 두 전이
$3\!\to\!2\mathrm L$ 와 $2\mathrm L\!\to\!2$ 는 독립한 둘이 아니라, 사이에 낀 \emph{stage-2L}(부분 채운 중간 stage)을
각각 진입$\cdot$이탈하는 \emph{분리쌍}이다. stage-2L 은 부분 채움의 배열 자유도가 $-T\Delta S$ 로 자유에너지를
낮춰 안정화되는 상이라 대략 $\gtrsim10\,^\circ$C 에서 뚜렷하고, 그 아래에서는 안정성을 잃어 쌍이 하나의
$3\!\to\!2$ 처리로 합쳐진다. 두 분리 전이의 반응 엔트로피가 갈리는 것이 관측 서명이다 --- 식~\eqref{eq:Uj} 의
$\partial U_j/\partial T=\Delta S_{\rxn,j}/F$ 로, 두 중심의 \emph{분리 속도}는 엔트로피 차
$\Delta(\Delta S)\equiv\Delta S_{\rxn}^{3\to2\mathrm L}-\Delta S_{\rxn}^{2\mathrm L\to2}$ 가 정한다:
\begin{equation}
\frac{\partial}{\partial T}\big(U^{3\to2\mathrm L}-U^{2\mathrm L\to2}\big)
=\frac{\Delta(\Delta S)}{F},
\qquad \Delta(\Delta S)\approx29\ \mathrm{J/(mol\,K)}
\;\Rightarrow\;
\frac{29}{96485}\approx0.30\ \mathrm{mV/^\circ C}.
\label{eq:gr2l-split}
\end{equation}
이 $\Delta(\Delta S)\approx29$ J/(mol\,K) 는 분리쌍 값 $\Delta S_{\rxn}^{3\to2\mathrm L}\approx+15$(고V 쪽)$\cdot$
$\Delta S_{\rxn}^{2\mathrm L\to2}\approx-14$(저V 쪽)의 \emph{차}이며(진입 시 배열 자유도 증가$\cdot$이탈 시 감소로
부호 반전), 우리 온도 스윕 재현이 $0.271\ \mathrm{mV/^\circ C}$ 로 이 예측의 $\sim90\%$ 를 되돌린다. 이때
$\Delta(\Delta S)$ 는 표~\ref{tab:staging} 초기값($3\!\to\!2\mathrm L{:}\,0,\ 2\mathrm L\!\to\!2{:}\,-5$, 차 $=5$)이
아니라 그 초기값을 \emph{피팅으로 갱신}하는 대상이다(표 수치 임의 변경 없이, 분리 서명을 self-test 로 맞춤).

\begin{signbox}
\textbf{관측 귀결(부호).} 분리 속도 $+0.30\ \mathrm{mV/^\circ C}$ 는 온도가 오를수록 두 중심이 \emph{벌어짐}을
뜻한다. 간격이 폭 $w_j$ 를 넘으면 dQ/dV 에 두 봉우리가 분해되고, 못 넘으면 한 봉우리로 겹친다 --- 따라서
고온($45\,^\circ$C)에서 2-peak 로 분해되고 상온($25\,^\circ$C)에서 병합돼 보이며, stage-2L 자체가 불안정해지는
$\sim10\,^\circ$C 아래에서는 분리쌍이 사라진다. 이 온도 의존 stage-2L 은 조작 곡선맞춤이 아니라 독립 실측을
가진다 --- 탈리튬 방향 operando XRD 가 $43\,^\circ$C 에서 stage-2L 형성을 뚜렷이, $0\,^\circ$C 에서는 부재로
보고한다\cite{schmitt2022}(우리 재현 $0.271\ \mathrm{mV/^\circ C}$ 와 부호$\cdot$경향 정합). 흑연 삽입
엔트로피의 부호$\cdot$스케일이 독립 실측과 맞는 것도 이를 받친다\cite{reynier2003}.
\end{signbox}

\textbf{(c) 박스 --- stage-2L 분리쌍의 온도 지도.} 위를 한 식으로 닫으면 두 분리 전이의 중심이 온도를 따라
서로 반대로 움직이며 그 간격이 $\Delta(\Delta S)/F$ 로 벌어진다:
\begin{equation}
\boxed{\;
\begin{aligned}
U^{3\to2\mathrm L}(T)&=U^{3\to2\mathrm L}_{\mathrm{ref}}+\frac{\Delta S_{\rxn}^{3\to2\mathrm L}}{F}\,(T-T_{\mathrm{ref}}),\\
U^{2\mathrm L\to2}(T)&=U^{2\mathrm L\to2}_{\mathrm{ref}}+\frac{\Delta S_{\rxn}^{2\mathrm L\to2}}{F}\,(T-T_{\mathrm{ref}}),\\
\big|\,U^{3\to2\mathrm L}-U^{2\mathrm L\to2}\,\big|&\ \text{가 FWHM}\approx3.53\,w_j\ \text{를 넘는 }T\text{ 에서만 2-peak 분해}\quad(\Delta(\Delta S)\approx29\ \mathrm{J/(mol\,K)}).
\end{aligned}\;}
\label{eq:gr2l-box}
\end{equation}
이는 식~\eqref{eq:Uj} 의 $U_j(T)$ 를 분리쌍 두 전이에 각각 적용한 것뿐이며(그림~\ref{fig:UjT} 의 기울기 부호
지도와 같은 대수), 새 항이 없다.

\textbf{(d) 추가 후보 --- XRD 5-feature 세분 시드(선택, 기준 미변경).} 위 판정자$\cdot$온도 서명은
표~\ref{tab:staging} 의 \emph{4-전이 기준을 그대로} 쓴다. 그와 별개로, XRD 상도표는 리튬화 좌표를 따라 최대
다섯 개의 staging 특징(dilute $1'\!\to\!4$, $4\!\to\!3$, $3\!\to\!2\mathrm L$, $2\mathrm L\!\to\!2$, $2\!\to\!1$)까지
분해한다\cite{dahn1991}. 이 세분을 원하는 경우를 위해 \emph{선택적} 5-feature 시드
(\S\ref{sec:sum-staging} 4-전이 기준과 병존)를 두되, 다음을 명시한다:
(i) 이는 기존 4-전이 기준을 \emph{대체하지 않는} 추가 시드이고(기본 경로 bit-exact 보존),
(ii) 이 시드가 세분하는 것은 staging 특징의 \emph{개수$\cdot$위치}(XRD 구조 해상도)뿐이며 전이별
두-상/고용체 \emph{상성격은 배정하지 않는다} --- 어느 특징이 두-상이고 어느 것이 고용체인지의 \emph{확정}은
세분 전과 똑같이 \S\ref{sec:broadening-class}(실측 plateau) $+$ 피팅된 $\Omega_j$ 소관이며(위 경계 가드$\cdot$검산과
같은 규정), 본 시드가 그 분류를 \emph{다시 여는} 것이 아니며,
(iii) \emph{명칭 대응}(혼동 금지): 피팅 기준선은 어디까지나 표~\ref{tab:staging} 의 \emph{네 행}
$\{4\!\to\!3,\ 3\!\to\!2\mathrm L,\ 2\mathrm L\!\to\!2,\ 2\!\to\!1\}$(무변경)이고, 별도
5-feature 시드는 그 \emph{네 전이 라벨을 계승하고} 고전위 끝에 dilute
$1'\!\to\!4$ 특징 \emph{하나만} 더 얹는다(5 $=$ 4 라벨 $+$ 추가 1). 다만 시드의 U-\emph{위치}는 표의 표시 위치를
그대로 쓰지 않고 \emph{XRD 해상도로 재유도}된다 --- XRD 사다리는 dilute$\approx$0.210, $4\!\to\!3\approx$0.170,
$3\!\to\!2\mathrm L\approx$0.132, $2\mathrm L\!\to\!2\approx$0.116, $2\!\to\!1\approx$0.085 V 로, 표 네 행의 표시
$U\,(0.210/0.140/0.120/0.085)$와 위치가 다르다(라벨은 계승, 위치는 재유도). 이 더한 dilute 특징은
\S\ref{sec:broadening-class} 이 \emph{고용체}로 분류$\cdot$명명한 dilute$\to$stage4 영역이다(표는 이를 별행 없이
고전위 끝에 접어 둔다). 곧 시드의 다섯 라벨(XRD 구조 해상도)을 표의 네 전이 \emph{피팅 기준선}
(무변경)과 혼동해선 안 되며, 시드 U-위치는 그 기준선이 아니라 XRD 사다리다.
\emph{6 개 이상의 별개 \textbf{물리 상(phase)}으로 쪼개는 것은 XRD 미지원 curve-fitting} 이라 폐기한다(회절 상 없이
봉우리만 얹으면 $R^2$ 는 오르나 대응 상이 없다). \textbf{단} \emph{gallery} 수(곡선표현의 로지스틱 개수)는 상 수와
다르다 --- MSMR 문헌\cite{verbrugge2017} 은 흑연을 \emph{6-gallery} 로 적합하되 U0 가 거의 같은 gallery
\emph{쌍}(한 물리 peak 을 폭 다른 두 로지스틱으로)으로 표현한다(gallery $\ne$ 상). 이 고분해능 곡선표현은 새 상을
posit 하지 않으므로 정당하며, \emph{선택} 시드(6-gallery, $\omega_j\cdot U0_j$
문헌 anchor)로 병존한다 --- 곧 해상도 사다리는 \{4-전이 기준(기본)$\cdot$5-feature XRD$\cdot$6-gallery MSMR
$\cdot$7-gallery(v1.0.25.2, 아래)\}이고, 어느 것도 4-전이 피팅 기준선을 바꾸지 않는다(전부 additive opt-in).

\textbf{사다리의 마지막 칸 --- 흑연 7-gallery(v1.0.25.2 opt-in).} 전이 수를 3 부터 8 까지 올리며
\emph{모수 보정 지표}(전이당 파라미터 수를 벌점으로 무는 정보 기준)로 재면, 흑연은 $N{=}7$ 에서 최소이고
$N{=}8$ 에서 다시 오른다 --- 곧 데이터 자신이 세분을 멈추는 지점이 있다(무한 세분 방향이 아니다).
그리고 $N$ 을 4 에서 7 로 올릴 때 \emph{검출 피크 수는 3 그대로}인데 중심이 $5$ mV 안에 겹치고 폭만
2 배 이상 갈리는 \emph{근축퇴 쌍}의 수가 늘어난다 --- 위 MSMR 6-gallery 의 쌍 관용과 같은 현상이며,
늘어난 전이가 \emph{새 위치의 상이 아니라 한 봉우리의 형상 분해}임을 그대로 보여준다.
\emph{이 성질은 커널 선택과 무관하다}: 두-상 정칙용액 커널로 같은 스윕을 해도 $N$ 증가에 따라 같은 방향으로
근축퇴 쌍이 늘어난다. 곧 세분은 어떤 커널을 쓰든 나타나는 \emph{곡선 표현 해상도}의 성질이지 특정 커널의
인공물이 아니다. 따라서 7-gallery 역시 새 상을 posit 하지 않으며, \textbf{XRD 상 수(흑연 staging 4)는
불변}이다(\S\ref{sec:broadening-class} 위임). \emph{정직 단서}: 본 스윕은 단일 셀$\cdot$비평형
pOCV(완화 없음) 위에서 얻었고, 시드 파라미터는 피팅 override 전제의 tier-C 값이다.

\textbf{판정자 $\Omega\lessgtr2RT$ 의 실측 대응 --- 흑연은 두-상이되 \emph{marginal} 이다(v1.0.25.2).}
같은 흑연 반쪽셀을 정칙용액 커널의 물리 4 전이로 적합하면 전이별 $\Omega/RT$ 가 $1.9$--$2.5$ 구간에
\emph{모두} 모인다 --- 곧 판정자~\eqref{eq:gr2l-disc} 의 문턱 $2RT$ 바로 위아래다. 이는 독립 문헌 앵커 둘과
정합한다: 2-sublattice 정칙용액 적합의 면내 항 $\Omega_a\!\approx\!3.4\,RT$, 그리고 평균장 정칙용액 적합의
$\Omega_a\!\approx\!2.5\,RT$ --- 셋 다 $2RT$ 를 넘되 크게 넘지 않는다. 곧 \textbf{흑연 staging 전이의 두-상
성격은 강한 상분리가 아니라 문턱 근방의 \emph{약한} 상분리}이며, 이것이 평형 등온선이 수학적 델타가 아니라
매우 좁은 유한 폭으로 관측되는 물리적 이유다(\S\ref{sec:width} 이중지위 둘째 지위의 정량 배경).
\emph{두 명제를 분리해 읽을 것}: ``흑연 전이가 두-상이다''는 위 $\Omega$ 가 지지하지만, ``그러므로 두-상
커널이 곡선 표현에서 gallery 를 대체한다''는 별개 물음이고 후자는 성립하지 않는다(\S\ref{sec:inputs}).
\emph{정직 단서}: 탐색 경계에 붙은 $\Omega$ 는 식별된 값이 아니므로 물리량으로 읽지 말 것 --- 그 경우
$\Omega$ 는 두-상 파라미터가 아니라 봉우리를 좁히는 손잡이로 쓰이고 있다.

\textbf{가장 좁은 봉우리는 계측 분해 한계에 걸려 있다(v1.0.25.2).} 흑연 최저 전위 평탄($\sim$0.10 V)은
원자료에서 전위 폭 $0.6$ mV 남짓 안에 수천 점$\cdot$유의한 전하가 몰려 있어 반높이 폭이 $1$ mV 이하 --- 열
척도 $RT/F\!\approx\!25.7$ mV 의 \emph{약 1/25} 다. 대칭 로지스틱으로 이를 맞추려면 $n_j$ 를 $0.02$ 수준까지
내려야 하는데, 이는 고용체 해석으로는 나올 수 없는 폭이다(고용체의 하한이 곧 $n_j{\sim}1$). 두-상 평탄으로
읽으면 자연스럽다 --- 즉 이 봉우리는 위 $\Omega$ 판정과 같은 방향을 가리킨다. \emph{단} 완화 없는 pOCV 의
비평형 잔여가 평탄을 인위적으로 좁혔을 가능성이 배제되지 않았다 --- 이 갈림은 완화를 포함한 평형 프로토콜
데이터로만 판정된다(미실행).

\textbf{사다리의 v1.0.25 확장과 그 실측 근거 --- gallery $\ne$ 상의 재확인.} 사다리에 Si 쪽 최고 해상도 옵션
하나가 더해진다: \emph{7-gallery}(전이 중심 anchor 에 $0.433\cdot0.456$ V 를 포함, 순수 로지스틱, opt-in).
근거는 프로토콜 의존성이다 --- 완화 구간을 포함한 pOCV(p-ocvhold) 프로토콜의 Si dQ/dV 에서는 $0.43$--$0.46$ V
에 $c$-Li$_{15}$Si$_4$ 결정화 쌍에 대응하는 두 feature 가 뚜렷이 잡히지만, 완화 없는 pOCV(p-ocv)에서는
\emph{나타나지 않는다}. 곧 관측되는 feature 개수 자체가 측정 조건의 함수다. 흑연에서도 같은 성질이 확인된다 ---
같은 셀$\cdot$같은 온도의 dQ/dV 에서 봉우리 검출 문턱(prominence)을 낮추면 검출 봉우리 수가 $3\to4\to5$ 로
갈라지고, 고해상 gallery 로 적합했을 때 나오는 중심 $U$ 들은 새 위치가 아니라 기존 세 봉우리
위치($\approx105\cdot141\cdot227$ mV)에 $\Delta U<12$ mV 의 \emph{근축퇴 쌍}으로 붙는다 --- 위 MSMR
6-gallery 관용\cite{verbrugge2017}(한 물리 peak $=$ 폭 다른 두 로지스틱)과 정확히 같은 양상이다.
\emph{따라서 사다리를 올린다고 상이 늘지 않는다} --- 흑연의 XRD staging 전이 수(넷)와 물리 두-상
수(둘: $2\mathrm L\!\to\!2\cdot2\!\to\!1$)는 어느 해상도에서도 \emph{불변}이며(상성격 판정은
\S\ref{sec:broadening-class} 소관), 사다리가 바꾸는 것은 곡선 표현의 분해능뿐이다. 같은 이유로 gallery 세분은
skew $\alpha_j$(\S\ref{sec:eqpeak} 식~\eqref{eq:skewpeak})와 서로 \emph{대안}이며 동시 자유화 시 축퇴한다
(\S\ref{sec:inputs} 식별 가드). 봉우리 수를 상 수의 증거로 제시하지 말 것 --- 그 증거는 회절이 준다.

\medskip
\textbf{stage-2L 의 실재성 --- 실측 근거·반증 가능 예측·상한 규율.} stage-2L 이 곡선맞춤용 여분 상이 아니라 실재 상임은 세 가지가 뒷받침한다. \emph{첫째, 실측 근거.}
stage-2L 은 dQ/dV 를 맞추려 도입한 여분 봉우리가 아니라 Dahn 상도표\cite{dahn1991} 가 회절로 고정한 stage
이고, 탈리튬 방향에서도 온도 의존 stage-2L 이 operando XRD 로 독립 관측된다\cite{schmitt2022}. \emph{둘째,
반증 가능 예측.} 식~\eqref{eq:gr2l-split} 의 $0.30\ \mathrm{mV/^\circ C}$ 분리와 $\sim10\,^\circ$C 병합은
\emph{자유 파라미터가 아니라} $\Delta(\Delta S)$ 한 값이 정하는 온도 서명이라, 다온도 데이터가 이를 배반하면
가정이 깨진다(재현 $0.271$ 은 정합). \emph{셋째, 상한의 규율.} 세분 시드조차 다섯 특징이 XRD 상한이며, 6+ 로의
분할은 폐기한다(위 (d)).

\begin{warnbox}
\textbf{적용 한계.} (i) \emph{절대 $\Delta S$ vs 엔트로피 \emph{차}.} 분리를 지는 양은 쌍의 \emph{차}
$\Delta(\Delta S)\approx29\ \mathrm{J/(mol\,K)}$ 이지 개별 절대값이 아니다. $+15/-14$ 는 이 차를 내는 \emph{한}
배치이며 표~\ref{tab:staging} 초기값($0/-5$, 차 $5$)은 관측 분해를 못 내는 거친 출발점이라, 분리 서명은 그
초기값을 $\Delta(\Delta S)$ 축으로 갱신하는 \emph{피팅 대상}이다(표 미편집). (ii) \emph{온도 규모의 등급.}
$\sim10\,^\circ$C 병합과 $25/45\,^\circ$C 분해 여부는 \emph{다온도 데이터 없이는 미검증}인 정성$\cdot$자릿수
예측이다 --- 회사 다온도 셀 데이터가 들어와야 $0.30\ \mathrm{mV/^\circ C}$ 와 병합점이 확정된다.
\cite{schmitt2022} 는 부호$\cdot$경향의 독립 근거이나 본 문건의 정량 $0.30$ 을 그 논문이 직접 보고한 값은
아니다(tier B). (iii) \emph{분류는 §7 소관.} 본 소절은 두-상/고용체 분류를 \emph{다시 여는} 것이 아니라
\S\ref{sec:broadening-class} 의 실측-plateau 분류를 판정식으로 닫을 뿐이다 --- regsol2 $\Omega_j/RT$ 진단은 우리
피팅 산출(표~\ref{tab:staging} 의 $\Omega$ 열과 같은 tier)이지 전이별 문헌 anchor 가 아니다. (iv) \emph{5-feature
는 추가 후보.} 세분 시드는 선택적 병존 도구이며 기준 4-전이를 대체하지 않고, 전이별 상성격도 배정하지 않는다
(그 성격은 \S\ref{sec:broadening-class}$+$피팅된 $\Omega_j$). (v) \textbf{상온 단일 곡선
이득이 아니다.} stage-2L 세분을 \emph{상온 단일 등온선}에 더하면 $R^2$ 는 오히려 $-0.44\%$p \emph{하락}한다
($25\,^\circ$C 에선 두 peak 가 병합이라 세분이 여분 자유도로 작용). 곧 stage-2L 은 \emph{상온 커브 피팅 이득이
아니라 다온도 매트릭스에서만 드러나는 서명}이며, 그 지지 근거는 $\Omega>2RT$ 판정 $+$ T-split 재현($0.271$
mV/$^\circ$C) $+$ LCO 다-feature $+1.90\%$p 이지 상온 $R^2$ 가 아니다 --- 상온 $R^2$ 로 stage-2L 을 정당화하지 않는다.
\end{warnbox}

\begin{keybox}
\textbf{stage-2L 다리.} 판정자~\eqref{eq:gr2l-disc}($\dd\mu/\dd\theta|_{1/2}=4RT-2\Omega$)가 \S\ref{sec:width}
이중지위의 두 갈래를 한 부등호로 닫되, 전이별 두-상/고용체 \emph{확정}은 \S\ref{sec:broadening-class} $+$ 피팅된
$\Omega_j$ 소관이다(판정자는 분류를 다시 열지 않음). 가운데 분리쌍은 엔트로피 안정화 stage-2L 을 사이에 둔 한
쌍이며, 그 존재의 서명이 $\Delta(\Delta S)\approx29\ \mathrm{J/(mol\,K)}\Rightarrow0.30\ \mathrm{mV/^\circ C}$
온도 분리(식~\eqref{eq:gr2l-box}, 재현 $0.271$; operando XRD 부호 정합\cite{schmitt2022})라는 \emph{반증 가능}한
예측이다. 5-feature 세분은 tab:staging 4-전이 기준을 \emph{라벨 그대로} 보존하고 고전위 끝 dilute $1'\!\to\!4$
하나만 더 얹는(5 $=$ 4 $+$ 1) \emph{추가 후보}(선택 시드, 상성격은 \S\ref{sec:broadening-class} 소관)다. \emph{별개 물리 상}
6+ 는 XRD 미지원으로 폐기하나, gallery$\ne$상이라 6-gallery MSMR 곡선표현(선택 병존)은
정당 --- 해상도 사다리 \{4-전이(기본)$\cdot$5-feature$\cdot$6-gallery\}는 전부 additive 이고 4-전이 피팅 기준선을 바꾸지 않는다.
\end{keybox}

% 자산(v1.0.24 @5 반영 — master 재조정판 · refine_b 위임-airtight화): [GR2L-1 판정자 dμ/dθ|½=4RT-2Ω eq:gr2l-disc(eq:gpp 대칭점)·기준식이지 판결 아님]
%   [GR2L-2 §7 분류 정합 verifybox(regsol2 Ω/RT — 분류는 실측 plateau 소관·과대주장 금지)]
%   [GR2L-3 stage-2L 엔트로피 분리 Δ(ΔS)=29→0.30 mV/℃ eq:gr2l-split·재현 0.271·+15/-14 피팅대상]
%   [GR2L-4 온도지도 boxed eq:gr2l-box·schmitt2022 operando XRD 부호 정합]
%   [GR2L-5 5-feature 세분=추가 후보(개수·위치만·상성격 배정 안 함·§7 소관·gate-7 5라벨↔4행 = 4+추가1 대응)] [GR2L-6 적대 3점+정직 한계 warnbox(−0.44%p 상온 caveat)]
